% Generated; evidence-dependent academic executive summary.
La ampliación compara modelos de reconstrucción mediante fidelidad física, después de examinar la calidad numérica de las referencias. La pregunta exige una misma observación, errores de velocidad y vorticidad, convergencia y costo a la tolerancia. El documento conserva cinco capítulos y desplaza el catálogo de evolución a los apéndices.

El banco principal incluye 1050 ajustes sobre siete referencias 128$^2$, cinco ruidos y tres semillas nuevas; sus 3150 evaluaciones Raw/H/H+B comparten esos ajustes. El reanálisis previo conserva 2808 condiciones y 8424 mediciones de costo. Los controles físicos alcanzan 512$^2$ en Navier–Stokes, 256$^2$ en viscoelasticidad y 256$^2$ en un intervalo viscoplástico corto.

\begin{enumerate}
\item \textbf{Fidelidad de simulación.} La resolución se juzga por observable y horizonte. Oldroyd-B puede tener velocidad casi invariable entre mallas y esfuerzo aún sensible; el refinamiento corto HB no valida automáticamente su evolución larga.
\item \textbf{Comparación del modelo.} Raw compara estimadores sin conocimiento físico añadido. Las diferencias se estratifican por contaminación y referencia, y los fallos permanecen en los denominadores.
\item \textbf{Aporte de restricciones.} H impone solenoidalidad; H+B limita banda física. Identity recibe las mismas transformaciones para distinguir qué mejora procede del estimador y qué mejora procede del prior.
\item \textbf{Costo útil.} CPU, RSS y trabajo dominante complementan el tiempo. La frontera de recursos sólo tiene sentido una vez exigida fidelidad; una observación sin método admisible conserva ese resultado negativo.
\end{enumerate}

La cobertura del banco, cuando basta alguna de las diez opciones fijas, es 3/105 con Raw, 16/105 con H y 105/105 con H+B. Se registran 0 fallos y 8 ajustes no convergentes. Estos conteos describen el banco y las cuatro puertas predeclaradas, no probabilidades poblacionales.

El prior tiene un costo físico medible: en Multiescala NS elimina 10.84~\% de vorticidad limpia en norma relativa. Cualquier reconstrucción limitada a esa banda conserva al menos ese error frente a la referencia completa, incluso sin ruido. Aceptar una salida con puerta del 25~\% no equivale a fidelidad del 10~\%.

\textbf{Guía de lectura.} El capítulo 2 define supuestos y observables; el 3 fija unidades de comparación y controles; el 4 responde las cuatro preguntas anteriores y examina sensibilidad. Los apéndices permiten inspeccionar evoluciones, resolución y resultados históricos completos sin diluir la comparación central.

El error de reconstrucción se evalúa contra la solución que generó la observación. Las diferencias contra una malla fina acotan qué interpretación física puede darse a ese error, pero no convierten una solución numérica en verdad experimental. El campo limpio no selecciona hiperparámetros; las sensibilidades y los pisos de sesgo se interpretan después del ajuste.
